% ============================================================
\section{From $\partial^2=0$ to the Carrier Taxonomy and Persistence Modes}
\label{sec:carrier-forcing}
% ============================================================

This chapter proves the structural theorem that the single identity
$\partial^2=0$ on the balanced complex~$K_4$, propagated through the defect
tower of the intrinsic $1{+}2{+}3$ decomposition, determines the common carrier
taxonomy and closure grammar. The defect tower first fixes the finite carrier
classes, the common closure grammar, and the admissible law-side channels on
those carriers. Once the carrier-specific endpoint-rigidity and representative
hypotheses are imposed, the corresponding canonical normal forms coincide with
the familiar Schr\"odinger, Klein--Gordon, Maxwell, and Einstein equations.
The claim carried forward to Part~II is stronger than four separate symmetry
recoveries: one selected bundle carries one common selected law, and the later
standard equations are its carrier-level realizations.
The references below are included only to situate those later identifications
in the standard mathematical-physics literature:
Schr\"odinger~\cite{Schrodinger1926,Sakurai2020},
Klein--Gordon~\cite{Klein1926,Gordon1926,PeskinSchroeder1995},
Maxwell~\cite{Maxwell1865,Jackson1999},
Einstein~\cite{Einstein1915,Wald1984}. For uniqueness of the Einstein tensor
in four dimensions, see Lovelock's theorem~\cite{Lovelock1971}; the framework
derives the same constraint ($b=-a/2$) internally via contracted Bianchi
from~$d^2=0$.

\subsection{Primitive Channels and Geometric Lift}

\begin{definition}[Visible persistence mode]
\label{def:visible-persistence-mode}\lean{SelectedBundle.FirstReading, SelectedBundle.CarrierClassification}
Fix a realized family satisfying the temporal realization principle on a
selected cut. A \emph{visible persistence mode} on that cut consists of:
\begin{enumerate}[label=(\roman*), leftmargin=2em, itemsep=0.2em]
\item a visible law class on the selected observed bundle;
\item a compatibility signature specifying what must be preserved by that law;
\item a source-admissibility profile specifying how exchange and returned
      residual may enter visibly;
\item a minimality condition asserting that no proper visible subcarrier still
      preserves the same compatibility signature and the same source profile on
      the realized temporal class.
\end{enumerate}
Such a mode is \emph{irreducible} if the minimality condition holds.
\end{definition}

\subsection{Defect-Tower Carrier Forcing}
\label{subsec:defect-tower-forcing}

The intrinsic \(1+2+3\) decomposition induces a finite carrier filtration on
the selected bundle:
\[
U_1
\subset
U_1\oplus U_2
\subset
U_1\oplus U_2\oplus U_3 = V_{\mathrm{slot}}.
\]
The direct channels are the successive bands of this three-step refinement,
while geometry is the symmetric-response lift of the last direct band rather
than a fourth primitive channel. Thus the carrier taxonomy is fixed before any
carrier-specific representative is chosen, and the later endpoint-rigidity
hypotheses identify representatives within already forced classes.
The role of \(D=4\) here is structural: Chapter~4 identifies the balanced
four-incidence as the canonical nondegenerate local carrier in the counting
model, and the explicit Part~II three-law hypothesis later gives that carrier
its physical-law reading.

\begin{theorem}[Defect-tower complete carrier forcing]
\label{thm:defect-tower-complete-carrier-forcing}\lean{SelectedBundle.Reduction.defectTowerCompleteCarrierForcing}
Assume the structural interface of Part~II, the canonical balanced
realization of Part~I with its intrinsic $1{+}2{+}3$ decomposition
(Theorem~\ref{thm:intrinsic-six-slot-decomposition}), and the temporal
realization on a faithful visible selected cut. Then the defect tower of the
intrinsic channel filtration forces exactly four intrinsic carrier classes.
Under the endpoint-rigidity and representative hypotheses used in the flagship
sections, their canonical normal forms are:
\begin{enumerate}[label=\textup{(\alph*)}, itemsep=0.3em]
\item \textbf{Phase carrier $\to$ Schr\"odinger equation.}
  The reversible recursion on the $U_1$ channel
  (Theorem~\ref{thm:reversible-transport-recursion}) produces a monodromy
  whose eigenvalues are constrained by the HFT-2 energy budget
  (Theorem~\ref{thm:HFT2-general}) to satisfy $|B|<2$
  (Proposition~\ref{prop:monodromy-2x2}). This forces complex eigenvalues,
  hence a canonical complex generator $J^2=-I$. Isotropic scalarization
  (Theorem~\ref{thm:isotropic-quadratic-transport}) then yields
  $i\partial_t\psi = (-\kappa\Delta + V)\psi$.
\item \textbf{Wave carrier $\to$ Klein--Gordon equation.}
  The HFT-1 defect (Theorem~\ref{thm:HFT1-general}) through the first
  correction band $U_2\oplus U_3$ generates the first defect correction to
  scalar observation. Scalarization on $U_2$ yields
  $(\partial_t^2 - c^2\Delta + m^2)\varphi = 0$.
\item \textbf{Gauge carrier $\to$ Maxwell equations.}
  The inherited $d^2=0$ on the $U_3$ channel restricts to de~Rham exactness
  on the gauge band. The Hodge solver
  (Theorem~\ref{thm:stencil-forcing-principle}) then yields
  $\delta F = J$, $dF = 0$.
\item \textbf{Geometric carrier $\to$ Einstein equation.}
  The symmetric-square lift $\mathrm{Sym}^2(U_3)$
  (Theorem~\ref{thm:six-slot-carrier-is-symmetric-square-of-u3}) carries an
  $O(n)$ decomposition. Isotropic symmetric-response scalarization
  (Corollary~\ref{cor:realized-isotropic-symmetric-response-scalarization})
  yields $\mathrm{div}(a\,\mathrm{Ric} + b\,R\,g) = (a/2+b)\,dR = 0$,
  forcing $b = -a/2$ and hence $G + \Lambda g = \kappa T$.
\end{enumerate}
The tower terminates in exactly three refinement steps (the three stages of
the $1{+}2{+}3$ filtration). There are exactly three primitive channels and
exactly one geometric lift. No carrier-level free parameter survives.
\end{theorem}

\begin{proof}
Part~(a): The recursion interface of Part~II
(Theorem~\ref{thm:reversible-transport-recursion}) on the $U_1$ channel
produces the monodromy $\mathcal M = \bigl(\begin{smallmatrix} B & -1 \\ 1 &
0 \end{smallmatrix}\bigr)$ with trace $B$. By Proposition~\ref{prop:monodromy-2x2},
the eigenvalues are $\lambda_\pm = (B \pm \sqrt{B^2 - 4})/2$. The HFT-2
energy budget (Theorem~\ref{thm:HFT2-general}) forces
$\sum_n \|x_n\|^2 < \infty$, which requires $|\lambda_\pm| \le 1$, hence
$|B| < 2$. This forces $B^2 - 4 < 0$, giving complex conjugate eigenvalues
on the unit circle and the canonical complex generator $J^2 = -I$. Isotropic
scalarization (Theorem~\ref{thm:isotropic-quadratic-transport}) on the
resulting complex line yields the Schr\"odinger form.

Part~(b): The HFT-1 identity (Theorem~\ref{thm:HFT1-general}) gives
$d_h^2 = -\delta_h$, where $\delta_h$ is the defect. The first correction
band $U_2$ inherits this defect structure. Scalarization under the isotropic
quadratic transport theorem yields a second-order equation. The mass term
arises from the constant part of the defect operator restricted to $U_2$.

Part~(c): On $U_3$, the full $d^2 = 0$ restricts to de~Rham exactness on
the gauge band. The solver layer produces the Green operator when
$H^k(D) = 0$ (exactness). The Hodge decomposition and the stencil-forcing
principle (Theorem~\ref{thm:stencil-forcing-principle}) yield $\delta F = J$
and $dF = 0$.

Part~(d): The symmetric-square lift
(Theorem~\ref{thm:six-slot-carrier-is-symmetric-square-of-u3}) on
$\mathrm{Sym}^2(U_3)$ inherits $d^2 = 0$ as the contracted Bianchi
identity. Isotropic symmetric-response scalarization
(Corollary~\ref{cor:realized-isotropic-symmetric-response-scalarization})
decomposes the divergence as
$\mathrm{div}(a\,\mathrm{Ric} + b\,R\,g) = (a/2 + b)\,dR$. Setting
this to zero forces $b = -a/2$, which is exactly the Einstein tensor
$G = \mathrm{Ric} - \frac{1}{2}R\,g$.
\end{proof}

\begin{remark}[Endpoint-rigidity hypotheses identify representatives]
\label{rem:endpoint-rigidity-packages-as-consequences}\lean{EndpointRigidity.CommonReduction.surface}
The carrier-specific endpoint-rigidity hypotheses of the flagship sections
identify canonical representatives for the carrier classes forced here. The
defect tower already determines the intrinsic carrier types and their common
closure grammar. When the local rigidity hypotheses are available, the
representative theorem is immediate; when they remain explicit, they serve
only to identify the endpoint normal form within an already forced class.
\end{remark}

\begin{remark}[Scope of the present carrier theorem]
\label{rem:still-missing-autonomous-normal-form-collapse}\lean{SelectedBundle.Reduction.defectTowerCompleteCarrierForcing, EndpointRigidity.CommonReduction.surface}
The present chapter proves the common carrier taxonomy and closure grammar. A
stronger theorem would remove the remaining endpoint hypotheses by showing
directly that, on an autonomous carrier satisfying locality, compatibility,
symmetry, and minimality, the only admissible local normal forms are the
flagship representative equations. That autonomous normal-form collapse is not
part of the present certified surface, so the endpoint hypotheses remain
explicit.
\end{remark}

\begin{corollary}[Single-source carrier taxonomy and closure grammar]
\label{cor:single-source}\lean{SelectedBundle.Reduction.singleSourceIdentification}
The common carrier taxonomy and closure grammar derive from the single
structural identity \(\partial^2=0\) on the canonical balanced clique complex,
together with the tower calculus of Part~II. Under the explicit
endpoint-rigidity and representative hypotheses stated in the flagship
sections, the canonical normal forms on those carriers coincide with the
familiar equations. No additional structural law appears at the common
law-side level.
\end{corollary}

\begin{proof}
The carrier taxonomy and closure grammar are forced by
Theorem~\ref{thm:defect-tower-complete-carrier-forcing}. The later endpoint
identifications require exactly the endpoint-rigidity and representative
hypotheses recorded in the flagship sections, as emphasized in
Remark~\ref{rem:still-missing-autonomous-normal-form-collapse}.
\end{proof}

\begin{definition}[Three primitive persistence channels and one intrinsic lift]
\label{def:five-persistence-modes}\lean{SelectedBundle.CanonicalGeneration}
Assume the setting of Definition~\ref{def:physical-realization-principle} on a
selected observed bundle whose local balanced channel algebra carries the
intrinsic \(1+2+3\) decomposition of
Corollary~\ref{cor:intrinsic-symmetry-forces-123}. The primitive visible
classification considered in this chapter consists of three direct persistence
channels together with one intrinsic lift:
\begin{enumerate}[label=(\roman*), leftmargin=2em, itemsep=0.2em]
\item \emph{local-state persistence:} the primitive \(1\)-channel, a
      reversible scalar mode whose
      compatibility signature is phase-compatible transport and whose visible
      source profile is source-free up to exchange or returned residual;
\item \emph{propagative persistence:} the primitive \(2\)-channel, a
      reversible scalar mode whose
      compatibility signature is second-order transport with no preferred
      direction and whose visible source profile is again source-free up to
      exchange or returned residual;
\item \emph{comparative persistence:} the primitive \(3\)-channel, a
      one-form comparison mode whose
      compatibility signature is exactness and local gauge covariance, so that
      only local comparison class is physical and visible source appears as an
      exchange current;
\item \emph{causal-structural persistence:} not a fourth primitive channel,
      but the canonical symmetric-response lift of the primitive \(3\)-channel
      on the time-realizing cut. Its compatibility signature is
      divergence-compatible response on the selected symmetric-square lift, and
      its visible source appears as a preceding exchange tensor.
\end{enumerate}
\end{definition}

\begin{theorem}[Primitive-channel and lift separation on the selected cut]
\label{thm:persistence-mode-separation}\lean{SelectedBundle.CanonicalGeneration.surface}
Assume temporal realization on a selected cut and the compatibility packages
described in Definition~\ref{def:five-persistence-modes}. Then the three
primitive persistence channels and the causal-structural lift of
Definition~\ref{def:five-persistence-modes} are pairwise distinct in role.
More precisely, no one of the three primitive channels can be identified with
another without destroying its compatibility signature, and the
causal-structural lift cannot be identified with its comparative base without
replacing a comparison-class law by a symmetric-response law with a different
source profile.
\end{theorem}

(Proof in Appendix~\ref{sec:appendix-carrier}.)

\begin{theorem}[Resolution-integrated balance closure is not a primitive channel or intrinsic lift]
\label{thm:resolution_integrated_balance_closure_not_intrinsic_mode}\lean{SelectedBundle.Reduction.surface}
Assume temporal realization on a selected cut and let a visible balance
hierarchy be obtained there by integrating unresolved exchange over a
resolution family, so that its recognizable hydrodynamic law appears only
after closure of that hierarchy. Then this hydrodynamic closure is neither an
additional primitive persistence channel nor a further intrinsic lift on the
selected cut. It is an effective closure read from the
resolution-integrated balance hierarchy.
\end{theorem}

(Proof in Appendix~\ref{sec:appendix-carrier}.)

\begin{remark}[Interpretation of the primitive channels and geometric lift]
\label{rem:persistence-mode-physical-reading}\lean{SelectedBundle.FirstReading, SelectedBundle.CanonicalGeneration}
These roles distinguish four structurally different ways in which coherence can
persist through realized time:
\begin{enumerate}[label=(\roman*), leftmargin=2em, itemsep=0.2em]
\item persistence of a local state on the primitive \(1\)-channel;
\item persistence by reversible propagation on the primitive \(2\)-channel;
\item persistence of local comparison on the primitive \(3\)-channel;
\item persistence of causal structure under symmetric response on the
      canonical lift of that primitive \(3\)-channel.
\end{enumerate}
The carrier language supplies minimal mathematical realizations of the three
primitive channels and of the geometric lift. Effective balance laws such as
hydrodynamics belong to a different layer: they are read only after
resolution integration and closure of an unresolved hierarchy.
\end{remark}

\subsection{Canonical Carrier Realizations}

\begin{definition}[Canonical realization of persistence modes on the selected bundle]
\label{def:canonical-sector-generation}\lean{SelectedBundle.CanonicalGeneration}
Assume the setting of Definition~\ref{def:physical-realization-principle} on a
selected observed bundle whose local balanced channel algebra carries the
intrinsic \(1+2+3\) decomposition of
Corollary~\ref{cor:intrinsic-symmetry-forces-123}. The bundle satisfies
\emph{canonical realization of persistence modes} if:
\begin{enumerate}[label=(\roman*), leftmargin=2em, itemsep=0.2em]
\item local-state persistence admits a minimal phase-generating realization,
      yielding the phase carrier class;
\item propagative persistence admits a minimal reversible second-order
      realization, yielding the wave carrier class;
\item comparative persistence admits a minimal exactness-compatible one-form
      realization, yielding the gauge carrier class;
\item the primitive \(3\)-channel admits a minimal divergence-compatible
      symmetric-response lift, yielding the geometric carrier class.
\end{enumerate}
This is not a primitive dictionary from the dimensions \(1+2+3\) to named
laws. It states that the visible carrier classes are the three canonical
direct realizations of the primitive channels together with the canonical
symmetric-response lift carried by the same intrinsic channel algebra.
\end{definition}

\begin{definition}[Resolution-integrated effective balance closure]
\label{def:resolution-integrated-effective-balance-closure}\lean{SelectedBundle.CarrierClassification, SelectedBundle.Reduction.surface}
A selected observed bundle carries a
\emph{resolution-integrated effective balance closure} if the bundle supports a
faithful balance hierarchy obtained by integrating unresolved exchange over a
resolution family, and the recognizable hydrodynamic law on that hierarchy is
read only as the effective constitutive closure of its exact low-order balance
content.
\end{definition}

\begin{definition}[Carrier-level classification of the persistence modes]
\label{def:carrier-level-classification}\lean{SelectedBundle.CarrierClassification}
A \emph{carrier-level classification} of the persistence modes on the selected
observed bundle consists in proving that each generated visible carrier carries
a unique
canonical representative of its distinguished minimal compatible closure class:
phase, wave, gauge, and geometric, together with proving that any
resolution-integrated effective balance hierarchy on that bundle carries a
canonical effective hydrodynamic closure normal form.
\end{definition}

\begin{theorem}[Persistence-mode realization reduction]
\label{thm:canonical-sector-generation-reduction}\lean{SelectedBundle.Reduction.surface}
Fix a realized family satisfying the temporal realization principle on a
selected cut. Assume:
\begin{enumerate}[label=(\roman*), leftmargin=2em, itemsep=0.2em]
\item the selected visible effective law on that cut is exactly the law
      obtained by eliminating a residual return field on the hidden/interface
      sector;
\item the four intrinsic flagship carriers together with the
      resolution-integrated hydrodynamic lane are faithfully visible on one selected
      observed bundle in the sense of
      Theorem~\ref{thm:common-selected-bundle-assembly};
\item the selected observed bundle satisfies canonical realization of
      persistence modes in the sense of
      Definition~\ref{def:canonical-sector-generation};
\item the same selected observed bundle carries a resolution-integrated
      effective balance closure in the sense of
      Definition~\ref{def:resolution-integrated-effective-balance-closure}.
\end{enumerate}
Then the phase, wave, gauge, and geometric carriers are the canonical
realizations of the three primitive persistence channels together with the
canonical geometric lift on one selected observed bundle, the kinetic lane is
the resolution-integrated effective balance
hierarchy on that same bundle, and the carrier-level next step is exactly the
classification of
Definition~\ref{def:carrier-level-classification}.
\end{theorem}

(Proof in Appendix~\ref{sec:appendix-carrier}.)

\begin{theorem}[Existence of the three primitive channels and the geometric lift on the selected bundle]
\label{thm:existence-of-five-persistence-modes}\lean{SelectedBundle.Reduction.surface}
Assume the hypotheses of
Theorem~\ref{thm:canonical-sector-generation-reduction}. Then each of the three
primitive persistence channels and the geometric lift of
Definition~\ref{def:five-persistence-modes} is realized on the selected
observed bundle. More precisely:
\begin{enumerate}[label=(\roman*), leftmargin=2em, itemsep=0.2em]
\item local-state persistence is realized by the phase carrier class;
\item propagative persistence is realized by the wave carrier class;
\item comparative persistence is realized by the gauge carrier class;
\item causal-structural persistence is realized by the geometric carrier class
      as the canonical symmetric-response lift of the primitive \(3\)-channel.
\end{enumerate}
\end{theorem}

(Proof in Appendix~\ref{sec:appendix-carrier}.)

\begin{theorem}[Mode classification for the intrinsic flagship family]
\label{thm:flagship-mode-classification}\lean{SelectedBundle.CarrierClassification.surface, SelectedBundle.Reduction.surface}
Assume the hypotheses of
Theorem~\ref{thm:canonical-sector-generation-reduction}. Then every intrinsic
flagship visible law on the selected observed bundle, namely the phase, wave,
gauge, and geometric lanes, belongs to exactly one of the three primitive
persistence channels or the geometric lift of
Definition~\ref{def:five-persistence-modes}.
\end{theorem}

(Proof in Appendix~\ref{sec:appendix-carrier}.)

\begin{corollary}[Scoped exhaustion of the intrinsic flagship primitive-channel family]
\label{cor:flagship-family-exhausts-five-persistence-modes}\lean{SelectedBundle.CarrierClassification.surface, SelectedBundle.Reduction.surface}
Assume the hypotheses of
Theorem~\ref{thm:canonical-sector-generation-reduction}. Then, within the
intrinsic flagship family on the selected observed bundle:
\begin{enumerate}[label=(\roman*), leftmargin=2em, itemsep=0.2em]
\item each intrinsic flagship carrier class realizes exactly one structural
      role from Definition~\ref{def:five-persistence-modes};
\item each of the three primitive persistence channels is realized by exactly
      one intrinsic flagship carrier class, and the geometric carrier is
      exactly the canonical symmetric-response lift of the primitive
      \(3\)-channel;
\item there is no fourth primitive persistence channel in the intrinsic
      flagship family, and no additional intrinsic lift beyond the geometric
      one.
\end{enumerate}
\end{corollary}

(Proof in Appendix~\ref{sec:appendix-carrier}.)

\begin{remark}[Hydrodynamic and thermodynamic closure are not additional primitive channels or intrinsic lifts]
\label{rem:thermodynamic-closure-not-sixth-mode}\lean{SelectedBundle.Reduction.surface, SelectedCut.effectiveThermodynamicClosure}
The corollary above concerns the intrinsic flagship family carried by the
selected observed bundle. It does not create an additional intrinsic mode for
hydrodynamics or thermodynamics. In the present chapter, hydrodynamics is read
as a resolution-integrated effective balance closure, while thermodynamic
irreversibility is read as a closure law on returned residual and unrecovered
stock. Neither is an additional primitive channel nor an additional intrinsic
lift beside the three primitive channels and geometric lift listed in
Definition~\ref{def:five-persistence-modes}.
\end{remark}

\subsubsection{\texorpdfstring{The \(D=2\) Spectral Seed}{The D=2 Spectral Seed}}
\label{subsec:d2-spectral-seed}
\label{rem:d2-spectral-seed}\lean{d2_spectral_seed_package}

By Corollary~\ref{cor:single-channel-unit-seed}, the unit-balanced seed admits
only one irreducible local channel. Therefore applying the temporal
realization bridge at that seed cannot produce scalar/vector/tensor branching:
the selected visible law there must remain one-fiber.
Theorem~\ref{thm:temporal-realization-forbids-primitive-internal-creation}
then rules out primitive source creation on that seed, so its leading visible
law is source-free unless an exchange remainder is induced by coupling or
hidden-sector return.

Combined with the reversible transport/recursion interface of
Section~\ref{subsec:reversible-transport-recursion} and the phase-carrier
rigidity machinery of Part~II, this identifies the natural seed visible law as
a reversible phase/recursion law. Corollary~\ref{cor:d2-repeat-no-new-carrier}
adds that repeat transport cycles on the seed contribute only powers of the
same one-channel monodromy and therefore do not create new visible persistence
modes. Corollary~\ref{cor:d2-primitive-support-profile},
Corollary~\ref{cor:primitive-support-prime-refinement}, and
Corollary~\ref{cor:primitive-support-continuum-spectral-defect} then identify
that same seed as a canonical primitive-support tower with prime-index maximal
cyclic refinement steps and a unique continuum spectral defect form.

\begin{figure}[ht]
\centering
\resizebox{0.8\textwidth}{!}{%
\begin{tikzpicture}[node distance=1.6cm]
  \node[ccbox, align=center, minimum width=3.0cm] (seed) at (0,1.8)
    {$D=2$ unit-balanced seed};
  \node[ccbox, align=center, minimum width=3.1cm] (channel) at (0,0.0)
    {one irreducible\\local channel};
  \node[ccbox, align=center, minimum width=3.2cm] (law) at (0,-1.9)
    {reversible phase/recursion\\visible law};
  \node[ccbox, align=center, minimum width=3.2cm] (support) at (0,-3.8)
    {primitive-support tower};
  \node[ccbox, align=center, minimum width=3.4cm] (continuum) at (0,-5.7)
    {unique continuum spectral\\defect form};

  \node[ccbox, align=center, minimum width=3.0cm] (repeat) at (4.8,-2.6)
    {repeat cycles contribute\\powers only};
  \node[ccbox, align=center, minimum width=3.0cm] (prime) at (4.8,-4.7)
    {maximal cyclic refinement\\is prime-index};

  \draw[ccarrow] (seed) -- (channel);
  \draw[ccarrow] (channel) -- (law);
  \draw[ccarrow] (law) -- (support);
  \draw[ccarrow] (support) -- (continuum);
  \draw[ccdarrow] (law.east) -- ++(0.55,0) |- (repeat.west);
  \draw[ccdarrow] (support.east) -- ++(0.55,0) |- (prime.west);
\end{tikzpicture}%
}
\caption{The \(D=2\) spectral seed. The one-channel seed forces a reversible
phase/recursion law; new spectral data arise only on primitive supports.}
\label{fig:d2-spectral-seed}
\end{figure}

This statement is operator-theoretic only: no arithmetic prime
identification is used. The \(D=2\) seed remains separate from the flagship
carrier class and furnishes the spectral stage beneath the nontrivial
\(D=4\) bundle.

Writing the forced continuum form as
\[
\mathcal E_{\infty}^{\mathrm{prim}},
\]
one natural external bridge problem is the identification
\[
\mathcal E_{\infty}^{\mathrm{prim}} = W_{\xi},
\]
where \(W_{\xi}\) denotes the classical completed Weil quadratic form on the
same log-time test class in Weil's positivity formulation of the explicit
formula \cite{Weil1952,Bombieri2000}. Any arithmetic realization of this
bridge would lie outside the operator calculus developed here. No such
identification is assumed here.

\subsubsection{Inherited Compatibility and Closure}

The flagship sections use these laws in five concrete representation
types. In each case, the carrier statement is proved inside the corresponding
section from the realized compatibility data, and the closure statement is the
inherited minimal-order closure mechanism already proved on the exported
structural interfaces.

\begin{remark}[Compatibility and order are inherited from the earlier structural interfaces]
\label{rem:compatibility-is-inherited}\lean{ObserverSelection.InheritedCompatibility}
No new compatibility identity is introduced at this stage. The compatibility
identities are already present in the exported theorem interfaces of
the earlier structural development and are only filtered by the selected
observer cut. Depending on
the realized family, these already include:
\begin{enumerate}[label=(\roman*), leftmargin=2em, itemsep=0.2em]
\item characteristic cut conservation and characteristic boundary reduction, by
      Theorem~\ref{thm:cut-conservation} and
      Corollary~\ref{cor:characteristic-boundary-reduction};
\item observer-side and dynamics-side triad compatibility, by
      Proposition~\ref{prop:observer-triad} and
      Theorem~\ref{thm:dynamics-triad};
\item reversible recursion and second-order wave transport on selected
      carriers, by
      Section~\ref{subsec:reversible-transport-recursion} together with the
      local datum-forcing interfaces cited in the wave and phase flagships;
\item nilpotent complex compatibility \(D_{k+1}D_k=0\) and the exactness/solver
      layer on finite Hilbert complexes, by
      Definition~\ref{def:hilbert-complex},
      Definition~\ref{def:exactness}, and
      Definition~\ref{def:solver-exactness};
\item local gauge covariance and unique gauge-compatible local datum on
      selected channels, by
      Corollary~\ref{cor:realized-local-gauge-covariance} and
      Theorem~\ref{thm:gauge-compatible-local-datum-forcing};
\item intrinsic symmetry scalarization on selected carriers, by
      Corollary~\ref{cor:realized-self-adjoint-scalarization} and
      Corollary~\ref{cor:realized-isotropic-symmetric-response-scalarization};
\item distinguished admissible minimal-order classes and their realized
      forcing, by
      Theorem~\ref{thm:minimal-admissible-transport-order},
      Corollary~\ref{cor:minimal-transport-order-controls-jet}, and
      Corollary~\ref{cor:realized-transport-order-forcing}; whenever a local
      representative family is later exhibited, Chapter~9 then verifies that
      representative by exact-core realization or constructive continuum
      forcing.
\end{enumerate}
Accordingly the later carrier sections ask only two questions: which inherited
compatibility identities are faithfully visible on the selected cut, and which
inherited minimal admissible closures preserve them there? Temporal
realization contributes the observer-side selection of the cut together with
the restriction that visible productions enter only as exact balanced
exchange.
\end{remark}

\subsubsection{Least-Motion Realization Data}

\begin{definition}[Characteristic observer datum]
\label{def:characteristic-observer-datum}\lean{ObserverSelection.CharacteristicDatum}
Fix a horizon \(h\), a characteristic horizon realization \(R\), and a
selection datum \(\mathsf S_h(R)\) in the sense of
Definition~\ref{def:selection-datum}. A \emph{characteristic observer datum}
at scale \(h\) consists of:
\begin{enumerate}[label=(\roman*), leftmargin=2em, itemsep=0.2em]
\item the selected horizon projector \(P_h\) and visible carrier
      \[
      \mathcal X_h:=P_h\cH^\bullet;
      \]
\item the lower-bound operator \(D_h\), candidate reduced-law family
      \(\mathcal Q_h(R)\), compatibility package, local-order filtration, and
      normalization package \(\mathcal N_h\) carried by \(\mathsf S_h(R)\);
\item the reconstruction and promotion data needed to place the selected family
      in the continuum-identification and admissible-promotion interfaces of
      the continuum bridge, together with the induced continuum carrier
      \(\mathcal X_{\chi^\ast}\) and continuum law space
      \(\mathcal Y_{\chi^\ast}\).
\end{enumerate}
\end{definition}

\begin{definition}[Least-motion characteristic observer]
\label{def:least-motion-characteristic-observer}\lean{ObserverSelection.LeastMotion}
A characteristic observer datum at scale \(h\) is a
\emph{least-motion characteristic observer} if:
\begin{enumerate}[label=(\roman*), leftmargin=2em, itemsep=0.2em]
\item it is faithful and reconstruction-compatible in the sense required by
      the continuum bridge;
\item it is closure-admissible, so the reduced admissible set
      \[
      \{Q\in\mathcal Q_h(R): D_h+Q\succeq 0\}
      \]
      is nonempty and well posed;
\item no proper visible subcarrier of \(\mathcal X_h\) still satisfies
      faithful reconstruction, compatibility preservation, and
      closure-admissibility;
\item its selected law class is stable under the admissible-promotion map of
      the continuum bridge;
\item among observer data satisfying (i)--(iv), it minimizes observer motion,
      meaning the combined representational drift of the projection, the
      minimal admissible completion, and the promoted law across the chosen
      horizon family;
\item it is unique up to horizon-preserving equivalence.
\end{enumerate}
\end{definition}

\begin{remark}[Triad consequence]
\label{rem:triad-consequence-observer-selection}\lean{ObserverSelection.TriadConsequence}
When a least-motion characteristic observer exists and is unique, the
Lagrangian/Eulerian/characteristic triad is read as a consequence of the
selected cut:
\[
\begin{aligned}
\Pi_{\chi^\ast}(\text{bulk-following description})
&=
\Pi_{\chi^\ast}(\text{boundary-observing description})\\
&=
\Pi_{\chi^\ast}(\text{action/generator description}).
\end{aligned}
\]
The surviving symmetry is then recovered as the stabilizer of the minimizing
observer datum rather than inserted by hand. In this sense, observer
covariance is the abstract symmetry statement, and the triad is its first
realized instance.
\end{remark}

\subsection{Common Structural Consequences}

\begin{theorem}[Common selected-bundle assembly]
\label{thm:common-selected-bundle-assembly}\lean{SelectedBundle.Assembly.surface}
Fix a realized family satisfying the temporal realization principle on a
least-motion characteristic observer datum with selected visible effective law
\(\mathcal L_{\mathrm{vis},h}(z)\). Assume that the phase, wave, kinetic,
gauge, and geometric intrinsic carriers together with the hydrodynamic balance
carrier on the selected cut are all faithful,
reconstruction-compatible, and stable under admissible promotion in the sense
of Definition~\ref{def:characteristic-observer-datum}.
Then:
\begin{enumerate}[label=(\roman*), leftmargin=2em, itemsep=0.2em]
\item these carriers are associated visible carriers of one selected observed
      bundle
      \[
      \mathcal X_h=P_h\cH^\bullet;
      \]
\item each induced visible law is obtained from the same selected
      visible effective law by restriction to the corresponding carrier
      together with the inherited selected endpoint closure class of
      Definition~\ref{def:selected-visible-effective-law};
\item once a canonical representative has been fixed on each carrier, the
      distinction between the flagship laws enters only at the final
      carrier-level endpoint classification and notation.
\end{enumerate}
\end{theorem}

(Proof in Appendix~\ref{sec:appendix-carrier}.)

\begin{theorem}[Common endpoint-rigidity reduction on the selected cut]
\label{thm:common-endpoint-rigidity-reduction}\lean{EndpointRigidity.CommonReduction.surface}
Fix a realized family satisfying the temporal realization principle on a
least-motion characteristic observer and an associated visible carrier
\(\mathcal X_h^{\mathrm{vis}}\) on the selected cut whose intrinsic symmetry
decomposition lies in the scalarization surfaces already forced upstream:
every non-geometric irreducible channel is an observable-irreducible selected
channel in the sense of
Definition~\ref{def:observable-irreducible-realized-class}, while every
geometric symmetric-response sector is an isotropic orthogonally equivariant
symmetric-tensor channel in the sense of
Corollary~\ref{cor:realized-isotropic-symmetric-response-scalarization}.
Then the selected endpoint closure class on
\(\mathcal X_h^{\mathrm{vis}}\) already has a unique exact carrier-side
representative, hence is unique modulo continuum equivalence, and is determined
entirely by scalar data on the intrinsic symmetry channels of that carrier.
Consequently the remaining carrier-level step is only to identify a natural
representative of that constructively forced class.
\end{theorem}

(Proof in Appendix~\ref{sec:appendix-carrier}.)
